\section{Setup}


\subsection{The Problem}
We first define the problem rigorously and introduce essential concepts for the proof.

Assume a fixed plane $\RR^2$ with the $xy$-coordinate system, so the $x$ and $y$ coordinate each refers to the first (horizontal) and second (vertical) coordinate respectively.
First we set up the hallway.
\textcolor{blue}{[A figure of hallway will help a lot]}
\begin{definition}
The \emph{hallway} $L = I_H \cup I_V$ is the union of its \emph{horizontal side} $I_H = (-\infty, 1] \times [0, 1]$ and \emph{vertical side} $I_V = [0, 1] \times (-\infty, 1]$. 

Define the \emph{horizontal strip} $H = \RR \times [0, 1]$ and \emph{vertical strip} $V = [0, 1] \times \RR$. 
With this, $I_H = L \cap H$ and $I_V = L \cap V$.
\end{definition}
We want to move the sofa from the hallway's horizontal side $I_H$ to its vertical side $I_V$.
We also give names to the sides and vertices of the hallway $L$ for future use. The convention is that the sides are in lowercase and the vertices are in boldcase lowercase. The names of sides follow that of \cite{romik2018differential}.
\begin{definition}
Let $a$, $b$, $c$ and $d$ be the sides 
\begin{align*}
a & = \{1\} \times (-\infty, 1] \qquad b = \{0\} \times (-\infty, 0] \\
c & = (-\infty, 1] \times \{1\} \qquad d = (-\infty, 0] \times \{0\}
\end{align*}
of $L$. 
\end{definition}
\begin{definition}
Let $\xx = (0, 0)$ and $\yy = (1, 1)$ be the inner and outer corner of the hallway $L$ respectively.
\end{definition}

We will describe the movement by a continuum of orientation-preserving isometries of $\RR^2$.
First, we use the following notations for translation and rotation maps.
\begin{definition}
For any vector $v \in \RR^2$, let $T_v$ be the translation map $T_v(w) = w + v$ of the plane by $v$.
For any $v \in \RR^2$ and $C \subset \RR^2$, write $T_v(C)$ as $C + v$.
For any angle $\theta \in \RR / 2 \pi \ZZ$, let $\Rot_\theta$ be the rotation map of $\RR^2$ along the center of the plane with a counterclockwise angle $\theta$.
\end{definition}
We adopt the following equivalent definition of an orientation-preserving isometry of $\RR^2$. For the proof of equivalence to a usual definition, see \textcolor{blue}{[Insert a suitable geometry book here]}.
\begin{definition}
A map $f \colon \RR^2 \ra \RR^2$ is called a orientation-preserving isometry of $\RR^2$ if
it is the composition $T_v \circ \Rot_\theta$ of a rotation followed by a translation.
\end{definition}
The following theorem allows us to identify the space $\cT$ of all orientation-preserving isometries of $\RR^2$ with $\RR^2 \times \RR / 2 \pi \ZZ$.
\begin{theorem}
An orientation-preserving isometry of $\RR^2$ is represented uniquely as the composition $T_v \circ R_\theta$ of a rotation followed by a translation.
Consequently, this gives a bijection $\iota \colon \RR^2 \times \RR / 2 \pi \ZZ \ra \cT$ defined as $(v, \theta) \mapsto T_v \circ R_\theta$.
\end{theorem}
\begin{definition}
For any orientation-preserving isometry $f = T_v \circ R_\theta$ of $\RR^2$, call $-\theta$ as the rotation angle of $f$.
Note the minus sign in front of $\theta$. This is to measure the rotation positively in the \emph{clockwise} direction.

We give the space $\cI$ of all orientation-preserving isometries of $\RR^2$ the topology of $\RR^2 \times \RR / 2 \pi \ZZ$ mapped under the bijection $\iota$.
In particular, the rotation angle is a continuous map on $\cI$.
\end{definition}

Now we are ready to describe a sofa by its shape and movement.
\begin{definition}
\label{def:sofa}
A \emph{sofa} $(S, f_t)$ is a nonempty, compact, connected subset $S$ of $\RR^2$ called the \emph{shape} of sofa equipped with a \emph{movement}, which is a rigid motion that moves $S$ from $I_H$ to $I_V$ inside $L$.
Formally speaking, the movement of sofa $S$ is a sign-preserving isometry $f_t$ of $\mathbb{R}^2$ continuously parametrized by time $0 \leq t \leq T$ for some $T \geq 0$, so that it satisfies the following conditions.
\begin{enumerate}
	\item The moving sofa $S_t = f_t(S)$ at any time $0 \leq t \leq T$ is always a subset of $L$.
	\item (Initial condition) The sofa $S_0 = f_0(S)$ at the beginning of the movement is a subset of $I_H$.
	\item (Final condition) The sofa $S_T = f_T(S)$ at the end of the movement is a subset of $I_V$. 
\end{enumerate}
\end{definition}

Now the moving sofa problem can be formalized clearly.
\begin{problem}[Moving Sofa Problem]
Does the maximum value of the area of the shape of a sofa exist?
If so, which sofa attains such a maximum value?
\end{problem}
We aim to prove that the answer is attained by Gerver's sofa \cite{gerver1992moving}.
\begin{conjecture}
Gerver's sofa \cite{gerver1992moving} attains the maximum possible area of a sofa.
\end{conjecture}

\subsection{Rotation Angle}

The \emph{rotation angle} is an important constant that determines the nature of a sofa.
\begin{definition}
The \emph{rotation angle} $\beta$ of a sofa $(S, f_t)$ with time interval $t \in [0, T]$ is the signed angular difference of $f_0$ and $f_T$ that tracks the full $2\pi$ rotations of $f_t$ during the movement.

More precisely, for each time $t \in [0, T]$, let $\theta(t) \in \RR / 2 \pi \ZZ$ be the angle the orientation-preserving isometry $f_t$ rotates the plane clockwise.
Then the map $\theta \colon [0, T] \ra \RR / 2 \pi \ZZ$ is continuous by the topology of $\cT$.
Take any continuous lift $\tilde{\theta} \colon [0, T] \ra \RR$ of $\theta$ 
so that the map $\tilde{\theta}$ composed with the modulus map $\RR \ra \RR / 2 \pi \ZZ$ is equal to $\theta$.
Then the rotation angle $\beta$ of $(S, f_t)$ is defined as $\tilde{\theta}(T) - \tilde{\theta}(0)$, which is independent of the choice of $\tilde{\theta}$.
\end{definition}

It turns out that we only need to consider sofas of rotation angle in between 0 and $\pi / 2$ inclusive to find the maximum area.
\begin{lemma}
\label{lem:rot-angle}
Any sofa of rotation angle $\beta > \pi / 2$ admits a movement with rotation angle $\pi / 2$. 
Any sofa of rotation angle $\beta < 0$ has an area at most $\sqrt{2}$.
\end{lemma}
\begin{proof}
\textcolor{blue}{Chore: Clean up the argument with more clear/crisp description}

Take any sofa $S$ of rotation angle $\beta > \pi / 2$ with movement $f_t$ ($t \in [0, T]$). 
By intermediate value theorem, during the movement of $S$ there is a time $T' \in [0, T]$ where the moving sofa $S_{T'}$ is a clockwise rotation of the initial position $S_0$ by $\pi / 2$. 
Consequently, the isometry $f_{T'}$ rotates the plane clockwise by an angle of $\pi / 2$.
Applying $f_{T'}$ to the initial condition $S_0 \subset H$, we obtain $S_{T'} \subset f_{T'}(H)$ where $f_{T'}(H)$ is a translation of the vertical strip $V$.
This implies that $S_{T'}$ has a width of $\leq 1$ in the direction of $x$-axis.

Now consider the following movement of $S$.
Move $S$ from $S_0$ to $S_{T'}$ by following the movement $f_t$ from $t=0$ to $t=T'$.
Since  $S_{T'}$ is contained inside $L$, we can now translate $S_{T'}$ inside $L$ towards the positive $x$-axis direction until it touches the side $a$ of $L$.
Because $S_{T'}$ has a width of $\leq 1$ in the direction of $x$-axis, its translated copy touching side $a$ should be contained in $V$.
This produces a movement of $S$ with rotation angle $\pi / 2$.

Take any sofa of rotation angle $\beta < 0$. If $\beta \geq -\pi / 4$,
the condition $S_T \subset V$ implies that $S_0$ is a subset of both $H$ and a rotated copy of $V$ by an angle of $\beta$. 
The intersection of the strips form a parallelogram of area $\leq \sqrt{2}$, concluding the proof. 
For $\beta < -\pi / 4$, by the intermediate value theorem, a copy of $S$ rotated counterclockwise by $\pi / 4$ lies in $L$.
Looking at this in perspective of sofa, $S$ is inside a copy $L'$ of $L$ rotated clockwise by $\pi / 4$. 
The intersection of $H$ and $L'$ always have area $\sqrt{2}$, so again $S \subset J \cap L'$ has an area less than or equal to $\sqrt{2}$.
\end{proof}

\subsection{Moving Hallway}
An important observation used by previous works \cite{gerver1992moving, romik2018differential, kallus2018improved} is to look at the sofa movement in the frame of sofa and not the hallway.
That is, instead of fixing the hallway $L$ and moving the sofa $S_t$ inside $L$ over time $t$, 
we follow the frame of reference of the sofa, making $S$ stationary, and move the hallway $L_t$ so that $L_t$ always contains $S$.

\begin{definition}[Alternative definition of a sofa]
\label{def:moving-hallway}
A sofa $(S, f_t)$ consists of the followings.
First, a shape $S$ which is a nonempty, compact, connected subset of $\RR^2$.
Second, the \emph{hallway movement} $g_t = f_t^{-1}$ which is a sign-preserving isometry continuously parametrized over a time interval $t \in [0, T]$ so that it satisfies the followings.

\begin{enumerate}
	\item The moving hallway $L_t = g_t(L)$ at any time $t \in [0, T]$ contains $S$.
	\item At the beginning of the movement, the sofa $S$ is contained inside $g_0(H)$, so consequently, it is also contained in $g_0(I_H) = g_0(H) \cap L_0$, the horizontal side of $L_0$ by (1).
	\item At the end of the movement, the sofa $S$ is contained inside $g_T(V)$, so consequently, it is also contained in $g_T(I_V) = g_T(V) \cap L_T$, the vertical side of $L_T$ by (1). 
\end{enumerate}
\end{definition}

Without loss of generality, we can assume that $S$ is already positioned in $S_0$, so that $g_0$ is the identity map of $\RR^2$, and thus $S$ is contained in $H$ in particular. 
With such an assumption, the three conditions of Definition \ref{def:moving-hallway} can be represented as the following single statement.
\begin{equation}
\label{eq:intersection}
S \subseteq H \cap g_T(V) \cap \bigcap_{0 \leq t \leq T} L_t.
\end{equation}
In this way, the sofa $S$ is a connected subset of a connected component of the intersection in Equation \eqref{eq:inter}.
So the problem of maximizing the area of $S$ becomes the problem of maximizing the area of the connected component of the \emph{intersection} of two strips $H$, $g_T(V)$ and the rotating hallways $L_t$.

Define the sides and corners of the moving hallway as the following.
\begin{definition}
Given a fixed sofa $(S, f_t)$ with hallway movement $g_t = f_t^{-1}$, denote the sides and corners of $L_t = g_t(L)$ corresponding to the sides $a, b, c, d$ and corners $\xx, \yy$ as $a(t), b(t), c(t), d(t)$ and $\xx(t), \yy(t)$.
So for example, $a(t) = g_t(a)$ and $\xx(t) = g_t(\xx)$.
\end{definition}
Note that this notation is ambiguous if there are more than one sofa $(S, f_t)$ being mentioned. 
We will put much emphasis on the inner corner $\xx(t)$ as a curve parametrized by $t$ later.

\subsection{Planar Convex Bodies}

In this section, we develop the theory of planar convex bodies. 
We will build the space of 

Notion of pointed corners

\subsection{History}

% add figure with caption here
% known as Hammersley's sofa, can be moved around with a continuous mix of rotation and transition.

The area of an explicit sofa gives a lower bound of the maximum sofa area. 
Since the problem was introduced, sofas of area $2.2074\dots $ (by Hammersley \cite{hammersley1968enfeeblement}), $2.21563\dots$ (by C. Francis \cite{guy1977monthly}) and $2.215649\dots$ (by R. Guy \cite{guy1977monthly}) were successively constructed over time. 
In 1997, Gerver \cite{gerver1992moving} found the sofa of currently best known area $2.2195\dots$ by considerations of the local maximality of a sofa.
The boundary of Gerver's sofa consists of 3 line segments and 15 analytic curves, making its derivation complicated.
Romik \cite{romik2018differential} simplified the derivation of Gerver's sofa in his work on ambidextrous sofa.

\begin{comment}
\begin{figure}[hbt]
  \includegraphics[width=0.9\textwidth]{img/gerver.png}
  \caption{Gerver's sofa with 3 line segments (V, XVIII, XIII) and 15 analytic curves. Shamefully copied from Wikipedia (need to put credit).}
  \label{fig:gerver}
\end{figure}
\end{comment}

On the other hand, the current best known upper bound $2.37$ of sofa area is found by Kallus and Romik in 2017 \cite{kallus2018improved} by a computer-assisted proof.
It is conjectured that the Gerver's sofa indeed attains the maximum sofa area, a speculation held by authors worked on this problem \cite{gerver1992moving, romik2018differential} that is supported experimentally by gradient descent computation on the problem's polygonal variant \cite{gibbs2014computational}.
Thus, under the assumption that Gerver's sofa, the resolution of moving sofa problem boils down to improving the upper bound of possible sofa area.
